\documentclass[conference]{IEEEtran}
\IEEEoverridecommandlockouts
\usepackage{cite}
\usepackage{amsmath}
\usepackage{graphicx}
\usepackage{xcolor}
\usepackage{tikz}
\usetikzlibrary{arrows.meta}
\usepackage{stfloats}
\usepackage{microtype}

\begin{document}
\raggedbottom

\title{Emulating a CXL Shared-Memory KV Cache over\\ One-Sided RDMA for Disaggregated LLM Inference}

\author{\IEEEauthorblockN{Chirayu Choudhary\IEEEauthorrefmark{1}, Karthikeyan Vaidyanathan, Jerome Anand}
\IEEEauthorblockA{International Institute of Information Technology Bangalore, India\\
chirayu.choudhary@iiitb.ac.in, karthikeyan.vaidyanathan@iiitb.ac.in, jerome.anand@intel.com}
\thanks{\IEEEauthorrefmark{1}Student Author, Integrated MTech.}}

\maketitle

\begin{abstract}
Disaggregated LLM serving moves key/value (KV) cache blocks from prefill to decode workers. TraCT replaces the RDMA path with CXL shared memory and, since CXL lacks cross-node coherence and atomics, adds a software protocol: offset-addressed objects, a shared prefix index and a two-tier lock. CXL pools are still rare, so we ask whether that protocol can run over commodity one-sided RDMA, and at what cost. On an InfiniBand cluster we build a POSIX shared-memory baseline standing in for CXL, an RDMA WRITE-with-immediate baseline, and a TraCT-style layer over RDMA READ/WRITE. Shared memory moves 4--8\,MiB blocks 2.3--2.7$\times$ faster than RDMA. The ported protocol is correct across three nodes; its lock-and-index path costs about 23\,$\mu$s per block, while sharing one memory host between writer and reader costs more.
\end{abstract}

\section{Introduction}
Disaggregated serving separates compute-bound prefill from latency-bound decode~\cite{distserve}, so every request's KV cache, often hundreds of megabytes, must move between workers, usually over RDMA~\cite{mooncake}. TraCT~\cite{tract} removes that network hop: workers share a CXL Type-3 memory pool used both as the transfer substrate and as a rack-wide prefix-aware KV cache, cutting mean time-to-first-token by up to 9.8$\times$ against RDMA and DRAM-cache baselines. As CXL memory lacks cross-host coherence and atomics, TraCT addresses objects by offsets, publishes only root objects such as the prefix index, and serializes writers through a two-tier lock: a per-node local lock, then a global array of per-node Idle/Waiting/Locked bytes arbitrated by a lock-manager thread.

\textbf{Limitation of the state of the art.} Multi-host CXL memory is not yet widely deployed; most research clusters, ours included, have InfiniBand but no CXL, and TraCT's results do not separate what the hardware buys from what the software design buys.

\textbf{Key insight.} The protocol needs only byte-granular remote loads and stores, which one-sided RDMA READ and WRITE provide, as RDMA shared-memory systems such as FaRM~\cite{farm} exploit. Its core mechanisms can therefore run over RDMA and be measured on commodity hardware. We contribute shared-memory and RDMA KV-transfer baselines, a TraCT-style layer over RDMA, and measurements on an InfiniBand cluster.

\textbf{Limitations.} Our shared-memory baseline is a single-node \texttt{memcpy} and does not model CXL device latency, which exceeds local DRAM~\cite{pond}; buffers live in host memory with no GPU in the path; one memory host serves all peers; and we use a bump allocator instead of TraCT's chunk allocator.

\section{Design}
A KV block holds 64 tokens of $B$ bytes; we sweep $B$ from 64 to 512\,KiB (4--32\,MiB blocks) and use 256\,KiB blocks for protocol tests. The baselines replay 20 requests of 1500, 3000, 4500 and 6000 tokens (1180 blocks); the RDMA layer draws 20 lengths in the same range (1193 blocks). Fig.~\ref{fig:timing} traces one block through each path.

\textbf{(a) Shared-memory baseline.} Prefill creates a pre-faulted POSIX segment holding a header with a process-shared mutex, an open-addressed table and a data area; per block it allocates under the mutex, copies the payload, then sets \texttt{valid}. Decode polls the entry. No NIC is involved, as with CXL load/store.

\textbf{(b) RDMA baseline.} A reliable-connected queue pair over native InfiniBand, set up by a TCP exchange of QP number, LID, PSN, address and rkey. Each block is one \texttt{RDMA\_WRITE\_WITH\_IMM} into a 4\,GiB registered ring on the decoder, the immediate carrying the request index, after one untimed warm-up block.

\textbf{(c) TraCT-style layer over RDMA.} A host node registers one region laid out like the CXL device: header, global lock array (17 locks $\times$ 8 node bytes), a 4096-slot prefix index and the data area, all referenced by offset. The host issues no RDMA; it only runs the lock manager, which sweeps the array and grants Locked to one Waiting node per lock (steps 1--3). Lock 0 guards the allocator and locks 1--16 guard 16 table partitions; the allocator needs its own lock because nodes holding different stripes would otherwise race on the offset counter. Collisions probe linearly within a partition. The payload is written unlocked (step 5); the entry, carrying \texttt{valid}=1, goes out in one WRITE only after the payload write completes, which is TraCT's visibility point (step 7), so readers need no lock (steps 8--9). Since the lock reaches remote memory through two function pointers, we first verified it over plain memory: with 3 simulated nodes of 2 threads, all nodes contended and the lock changed node over 5000 times with no violation, whereas a deliberately broken manager caused 209.

\begin{figure*}[t]
\centering
% sequence-style timing diagrams for the three KV transfer paths
\tikzset{
  life/.style={draw=black!55, line width=0.5pt},
  hdr/.style={draw=black!70, fill=black!6, rounded corners=1pt, inner sep=2pt, align=center, font=\fontsize{6}{6.6}\selectfont\bfseries},
  lbl/.style={font=\fontsize{5.7}{6.3}\selectfont, inner sep=1pt, align=left},
  ctl/.style={-{Stealth[length=3pt,width=2.4pt]}, line width=0.5pt, draw=black!80},
  poll/.style={-{Stealth[length=3pt,width=2.4pt]}, line width=0.5pt, dashed, draw=black!70},
  data/.style={-{Stealth[length=4pt,width=3.4pt]}, line width=1.8pt, draw=blue!65!black},
  loopbox/.style={draw=black!45, dashed, line width=0.4pt, rounded corners=1pt},
  note/.style={draw=black!40, fill=yellow!12, font=\fontsize{5.5}{6.1}\selectfont, inner sep=1.5pt, align=left},
  sub/.style={font=\fontsize{5}{5.5}\selectfont},
}
\newcommand{\sn}[1]{\textbf{\color{blue!60!black}#1}\,}
\newcommand{\sepbar}{\,$|$\,}

\begin{minipage}[t]{0.285\textwidth}\centering
\begin{tikzpicture}[x=1cm,y=-0.9cm]
  \def\A{0.35}\def\B{3.05}\def\C{4.6}
  \node[hdr] at (\A,0) {Prefill};
  \node[hdr] at (\B,0) {SHM segment\\[-1pt]{\fontsize{5}{5}\selectfont\mdseries hdr\sepbar table\sepbar data}};
  \node[hdr] at (\C,0) {Decode};
  \foreach \x in {\A,\B,\C} \draw[life] (\x,0.32) -- (\x,5.3);
  \draw[ctl] (\A,0.72) -- node[lbl,above]{shm\_open, mmap, pre-fault} (\B,0.72);
  \draw[ctl] (\C,1.05) -- node[lbl,above]{attach} (\B,1.05);
  \draw[loopbox] (0.02,1.25) rectangle (4.93,5.15);
  \node[lbl,anchor=north west,font=\fontsize{5.4}{6}\selectfont\itshape] at (0.04,1.27) {for each KV block};
  \draw[ctl] (\A,1.8) -- node[lbl,above]{\sn{1}lock (process-shared mutex)} (\B,1.8);
  \draw[ctl] (\A,2.2) -- node[lbl,above]{\sn{2}probe; off $\mathrel{+}=$ size; unlock} (\B,2.2);
  \draw[data] (\A,2.65) -- node[lbl,above]{\sn{3}memcpy payload $\rightarrow$ data[off]} (\B,2.65);
  \draw[ctl] (\A,3.1) -- node[lbl,above]{\sn{4}entry\{h, off, len\}; valid=1} (\B,3.1);
  \draw[poll] (\C,3.55) -- node[lbl,above]{\sn{5}poll valid} (\B,3.55);
  \draw[poll] (\C,3.78) -- (\B,3.78);
  \draw[data] (\B,4.25) -- node[lbl,above]{\sn{6}read data[off]} (\C,4.25);
  \node[note] at (1.7,4.75) {load/store only, no NIC\\[-1pt](stands in for CXL)};
\end{tikzpicture}\\[-1pt]
{\fontsize{7}{8}\selectfont (a) SHM baseline, one node}
\end{minipage}\hfill
\begin{minipage}[t]{0.285\textwidth}\centering
\begin{tikzpicture}[x=1cm,y=-0.9cm]
  \def\A{0.4}\def\B{2.5}\def\C{4.55}
  \node[hdr] at (\A,0) {Prefill};
  \node[hdr] at (\B,0) {HCA / IB\\[-1pt]{\fontsize{5}{5}\selectfont\mdseries RC queue pair}};
  \node[hdr] at (\C,0) {Decode\\[-1pt]{\fontsize{5}{5}\selectfont\mdseries reg.\ KV cache}};
  \foreach \x in {\A,\B,\C} \draw[life] (\x,0.32) -- (\x,5.3);
  \draw[ctl,{Stealth[length=3pt,width=2.4pt]}-{Stealth[length=3pt,width=2.4pt]}] (\A,0.72) -- node[lbl,above,pos=0.5]{TCP: QPN, LID, PSN, addr, rkey} (\C,0.72);
  \node[lbl,fill=white] at (2.48,1.0) {QP INIT$\rightarrow$RTR$\rightarrow$RTS (both sides)};
  \draw[ctl] (\A,1.42) -- node[lbl,above,pos=0.5]{warm-up WRITE\_IMM (not timed)} (\C,1.42);
  \draw[loopbox] (0.02,1.6) rectangle (4.93,5.15);
  \node[lbl,anchor=north west,font=\fontsize{5.4}{6}\selectfont\itshape] at (0.04,1.62) {for each KV block};
  \draw[ctl] (\C,2.15) -- node[lbl,above]{\sn{1}post RECV} (\B,2.15);
  \draw[data] (\A,2.6) -- node[lbl,above]{\sn{2}WRITE\_IMM (imm)} (\B,2.6);
  \draw[data] (\B,3.05) -- node[lbl,above]{\sn{3}DMA $\rightarrow$ ring[slot]} (\C,3.05);
  \draw[ctl] (\B,3.5) -- node[lbl,above]{\sn{4}RECV CQE (imm)} (\C,3.5);
  \draw[ctl] (\B,3.95) -- node[lbl,above]{\sn{5}send CQE: stop} (\A,3.95);
  \node[note] at (2.48,4.62) {one-sided WRITE; the immediate\\[-1pt]is the only thing that wakes decode};
\end{tikzpicture}\\[-1pt]
{\fontsize{7}{8}\selectfont (b) RDMA baseline, one node}
\end{minipage}\hfill
\begin{minipage}[t]{0.41\textwidth}\centering
\begin{tikzpicture}[x=1cm,y=-0.9cm]
  \def\A{0.45}\def\B{4.05}\def\C{6.85}
  \node[hdr] at (\A,0) {Writer\\[-1pt]{\fontsize{5}{5}\selectfont\mdseries compute01}};
  \node[hdr] at (\B,0) {Host: one MR + lock manager\\[-1pt]{\fontsize{5}{5}\selectfont\mdseries hdr\sepbar lock[17][8]\sepbar table[4096]\sepbar data \quad (compute00)}};
  \node[hdr] at (\C,0) {Reader\\[-1pt]{\fontsize{5}{5}\selectfont\mdseries compute02}};
  \foreach \x in {\A,\B,\C} \draw[life] (\x,0.32) -- (\x,5.3);
  \draw[loopbox] (0.05,0.5) rectangle (7.2,5.15);
  \node[lbl,anchor=north west,font=\fontsize{5.4}{6}\selectfont\itshape] at (0.07,0.52) {for each KV block};
  \draw[ctl] (\A,1.05) -- node[lbl,above]{\sn{1}WRITE lock[0][n]=W (allocator lock)} (\B,1.05);
  \draw[ctl] (\B+0.05,1.13) .. controls (\B+0.6,1.13) and (\B+0.6,1.45) .. (\B+0.05,1.45);
  \node[lbl,anchor=west] at (\B+0.55,1.29) {\sn{2}manager: one W$\rightarrow$L};
  \draw[poll] (\A,1.7) -- node[lbl,above]{\sn{3}READ lock[0][n] until L} (\B,1.7);
  \draw[ctl] (\A,2.15) -- node[lbl,above]{\sn{4}READ + WRITE next\_off; WRITE I} (\B,2.15);
  \draw[data] (\A,2.6) -- node[lbl,above]{\sn{5}WRITE payload $\rightarrow$ data[off], no lock} (\B,2.6);
  \draw[poll] (\A,3.05) -- node[lbl,above]{\sn{6}stripe s: WRITE W, READ until L} (\B,3.05);
  \draw[ctl] (\A,3.5) -- node[lbl,above]{\sn{7}READ probe; WRITE entry, valid=1; I} (\B,3.5);
  \draw[poll] (\C,4.0) -- node[lbl,above]{\sn{8}READ entry until valid} (\B,4.0);
  \draw[data] (\B,4.45) -- node[lbl,above]{\sn{9}READ data[off]} (\C,4.45);
  \node[note,anchor=west] at (0.15,4.4) {offsets, not pointers;\\[-1pt]s = 1 + partition(hash)};
  \node[lbl,font=\fontsize{5.4}{6}\selectfont\itshape,anchor=east] at (6.78,4.9) {reader takes no lock};
\end{tikzpicture}\\[-1pt]
{\fontsize{7}{8}\selectfont (c) TraCT-style layer over RDMA, three nodes}
\end{minipage}

\vspace{-6pt}
\caption{Per-block timing, time flowing down; thick arrows carry payload, dashed ones poll. All arrows in (c) are one-sided RDMA except local step 2.}
\label{fig:timing}
\vspace{-11pt}
\end{figure*}

\begin{figure*}[b]
\vspace{-14pt}
\centering
\includegraphics[width=0.315\textwidth]{figs/fig_throughput.pdf}\hfill
\includegraphics[width=0.315\textwidth]{figs/fig_footprint.pdf}\hfill
\includegraphics[width=0.315\textwidth]{figs/fig_breakdown.pdf}\\[-2pt]
\makebox[0.325\textwidth]{\footnotesize (a) Throughput vs.\ block size}\hfill
\makebox[0.325\textwidth]{\footnotesize (b) 32\,MiB blocks as footprint grows}\hfill
\makebox[0.325\textwidth]{\footnotesize (c) Writer time per 256\,KiB block}
\vspace{-4pt}
\caption{Results. In (a) the dotted line is the raw verbs RDMA WRITE peak (5.75\,GB/s).}
\label{fig:results}
\vspace{-4pt}
\end{figure*}

\section{Evaluation}
\textbf{Testbed.}\footnote{Plotting/data-processing scripts received light AI (Claude) assistance; experimental design, execution, and analysis are the authors' own.} An 8-node cluster (20--40 cores, 128\,GiB per node, Mellanox InfiniBand, 4096-byte MTU, GCC 4.8.5, libibverbs). Paths (a) and (b) run on one node, (b) through the HCA in loopback; (c) puts host, writer and reader on three compute nodes. A verbs micro-benchmark on the same node gives 1.34\,$\mu$s for a 1-byte RDMA WRITE and 5.75\,GB/s at 4\,MiB.

\textbf{Shared memory vs.\ RDMA (Fig.~\ref{fig:results}a).} Shared memory reaches 11.6\,GB/s at 4\,MiB and 13.5--13.6\,GB/s at 8--16\,MiB, 2.3--2.7$\times$ the RDMA baseline, which holds 4.95--4.97\,GB/s (86\% of the verbs peak). The gap is the NIC hop itself, as TraCT argues. It then falls to 11.9\,GB/s at 24\,MiB and 6.7\,GB/s at 32\,MiB.

\textbf{Cost of the ported protocol (Fig.~\ref{fig:results}c).} On one node the writer spends 71.7\,$\mu$s per 256\,KiB block against 48.7\,$\mu$s for a bare RDMA WRITE. The 23\,$\mu$s difference is the metadata path, at least ten small one-sided operations at 1.3--1.5\,$\mu$s each plus the manager's sweep delay, and would be 1--3\% of a 4--8\,MiB block. Sharing the single host costs more: with a concurrent reader the writer takes 112.6\,$\mu$s on one node and 133.4\,$\mu$s across three. At 4--8\,MiB the three-node writer, again with a concurrent reader, sustains 3.7--3.8\,GB/s, 31--34\% slower than the single-node RDMA baseline, while the reader reads at 3.9\,GB/s; separating host sharing from cross-node effects at these sizes is ongoing.

\textbf{Correctness.} In every run the reader found all 1193 index entries with matching keys, including runs where it started first and waited on in-flight writes (mean 1.6--9.5\,ms per block) and a run with two writers publishing disjoint key spaces.

\textbf{Memory footprint (Fig.~\ref{fig:results}b).} With 32\,MiB blocks both shared memory (7.8$\rightarrow$6.3\,GB/s) and the RDMA layer (3.9$\rightarrow$1.1\,GB/s) slow down over about 40\,GB written despite pre-faulting, while the RDMA baseline, which recycles a 4\,GiB ring, stays flat. This points to host-memory footprint rather than transport; we have not yet isolated the cause.

\section{Conclusion and Future Work}
TraCT's core software protocol ports to one-sided RDMA: its lock-and-index path adds about 23\,$\mu$s per block, and the larger costs are host sharing and bandwidth, the latter being what CXL hardware provides. Next: (1) sustained multi-writer contention on more nodes, to find where the single-threaded manager saturates; (2) RDMA compare-and-swap, which CXL lacks, in place of the manager; (3) TraCT's chunk allocator with eviction; (4) huge pages and NUMA pinning to isolate the footprint slowdown; (5) a serving engine with GPUDirect.

\begin{thebibliography}{5}
\footnotesize
\bibitem{tract} D.~Yoon, Y.~Min, H.~Kim, S.~H.~Noh, and J.~Kim, ``TraCT: Disaggregated LLM serving with CXL shared memory KV cache at rack-scale,'' arXiv:2512.18194, 2025.
\bibitem{distserve} Y.~Zhong \emph{et al.}, ``DistServe: Disaggregating prefill and decoding for goodput-optimized large language model serving,'' in \emph{Proc. USENIX OSDI}, 2024.
\bibitem{mooncake} R.~Qin \emph{et al.}, ``Mooncake: Trading more storage for less computation, a KVCache-centric architecture for serving LLM chatbot,'' in \emph{Proc. USENIX FAST}, 2025.
\bibitem{farm} A.~Dragojevi\'c, D.~Narayanan, O.~Hodson, and M.~Castro, ``FaRM: Fast remote memory,'' in \emph{Proc. USENIX NSDI}, 2014.
\bibitem{pond} H.~Li \emph{et al.}, ``Pond: CXL-based memory pooling systems for cloud platforms,'' in \emph{Proc. ACM ASPLOS}, 2023.
\end{thebibliography}

\end{document}
